\documentclass[11pt]{article}
\usepackage[margin=1.1in]{geometry}
\usepackage{amsmath,amssymb}
\usepackage{booktabs}
\usepackage{graphicx}
\usepackage[colorlinks=true,linkcolor=blue,citecolor=blue,urlcolor=blue]{hyperref}
\usepackage[numbers,sort&compress]{natbib}
\usepackage{xcolor}

\newcommand{\todo}[1]{\textcolor{red}{[TODO: #1]}}
\newcommand{\dscode}{\textsc{deepseek-chat}}
\newcommand{\qwcode}{\textsc{qwen-plus}}

\title{A Lean-Style Verifier for Seiberg Dualities Turns Weak Language
Models into Auditable Repair Agents\\
{\large\todo{working title, iterate with author}}}
\author{Xingyang Yu}
\date{\today}

\begin{document}
\maketitle

\begin{abstract}
We present \textsc{QFTCert}, a symbolic verifier that certifies
consistency of proposed Seiberg-duality claims in four-dimensional
$\mathcal{N}=1$ gauge theories (anomaly matching, superpotential
$R$-charges, central charges, bounded chiral-ring consistency), and use
it as an interaction environment for language-model agents that must
\emph{repair} deliberately broken duality claims. On a fixed benchmark of
145 depth-one perturbed claims, under a pre-registered analysis protocol
frozen before data collection, we find: (i) verifier-gated iteration
lifts repair success over single-shot attempts on both models tested
($+8.3$\,pp on \dscode{}, $+7.1$\,pp on \qwcode{}, Holm-adjusted
$p<0.002$); (ii) the \emph{optimal verifier-exploitation strategy is
model-dependent, with confirmatory significance in opposite directions}:
for \dscode{}, feedback content adds nothing over content-free retry and
independent verifier-filtered resampling dominates a strategy portfolio
($-10.3$\,pp), while for \qwcode{}, category-level feedback is worth
$+8.7$\,pp, interpretable obligation identities alone are worth
$+6.4$\,pp over structurally identical masked feedback, and the
portfolio dominates resampling ($+14.7$\,pp); (iii) the invariant across
both regimes is the verifier's cheap, machine-checkable certificate,
which powers whichever search strategy wins. The entire confirmatory
phase cost roughly \$11 of public API credit; all code, fixtures, frozen
protocol, and per-round audit logs are released.
\todo{finalize abstract wording after prose pass}
\end{abstract}

\section{Introduction}
\todo{motivation: proof assistants ground LLM reasoning in mathematics;
physics lacks the analogous machine-checkable substrate. QFTCert as a
domain proof-of-concept: consistency CERTIFICATES for duality claims
(Lean-style, not Lean-equivalent — no proof terms; the verifier
auto-judges). Honest scope: certificates of consistency, not proofs of
duality. Contributions list mirroring abstract claims (i)-(iii) +
benchmark + open release.}

\section{Consistency certificates for Seiberg duality}
\todo{summarize the verifier: pure-quiver claims, obligations
(anomaly matching, W R-charge = 2, ABJ anomaly freedom, central-charge
matching, R-graded bounded chiral-ring consistency at length L),
judge modes, out-of-scope semantics. Point to appendix for full
obligation list. Emphasize L=3 feedback / L=5 final anti-gaming pair.}

\section{Benchmark: perturbed duals at fixed depth}
\todo{seed families (dP0, dP1, dP2, F0, SPP, C3/Z2xZ2 + nodes),
mutation engine for ground-truth duals, four perturbation classes
(drop\_w\_term 38, flip\_w\_sign 13, r\_charge\_perturb 52,
rank\_perturb 42; n=145 repairable), deterministic regeneration
(seed 20260715), anti-memorization rationale, copy-cheat guard
(verbatim Theory-A copy trivially certifies; guard fired in live runs).}

\section{Experimental design and frozen protocol}
\todo{arms: single-shot, generic retry (K=5), verifier feedback
medium (K=5), masked feedback (E5 control, neutral positional
placeholders), best-of-11 (E4 control, independent draws, stop at first
final certificate); full-theory edit mode; complete component collection
+ deterministic stop-first portfolio replay (protocol section 4).
Endpoints E1--E5; paper-wide Holm over \{E1,E2,E4\}x\{2 models\};
logit-link GEE, fixture-clustered sandwich, standardized marginal risk
difference. Protocol frozen at commit 813fdb0 BEFORE any confirmatory
response; execution manifest with prompt/config/fixture hashes.
R=3 replications per cell. Disclosure: the class-level ``two-regime''
hypothesis from the pilot was post hoc and is NOT a confirmatory claim.}

\section{Pilot phase (exploratory)}
\todo{auditable qwen pilot (ss 13/145, gr 19, vf 21, det 32; paired
tables) in one subsection; UNVERIFIABLE legacy deepseek aggregates
(records lost) in a separate subsection, labeled as such, excluded from
inference. Frozen pilot top-line sentence goes here verbatim. The pilot
motivated: E5 masked control, E4 call-cap control, invalid-decomposition
reporting.}

\section{Confirmatory results}

Table~\ref{tab:primary} reports the frozen endpoints; Table~\ref{tab:perrep}
the underlying per-replication success counts.

\begin{table}[t]
\centering
\begin{tabular}{llrrr}
\toprule
model & endpoint & RD (pp) & 95\% CI & $p_{\mathrm{Holm}}$ \\
\midrule
\textsc{deepseek-chat} & E2 (\texttt{gr} $-$ \texttt{ss}) & +8.3 & [+3.8, +12.8] & 0.0015 \\
\textsc{deepseek-chat} & E1 (\texttt{vf} $-$ \texttt{gr}) & -1.8 & [-6.6, +3.0] & 0.4528 \\
\textsc{deepseek-chat} & E4 (portfolio $-$ \texttt{best\_of\_n}) & -10.3 & [-16.2, -4.5] & 0.0016 \\
\textsc{qwen-plus} & E2 (\texttt{gr} $-$ \texttt{ss}) & +7.1 & [+3.2, +11.1] & 0.0016 \\
\textsc{qwen-plus} & E1 (\texttt{vf} $-$ \texttt{gr}) & +8.7 & [+3.8, +13.7] & 0.0016 \\
\textsc{qwen-plus} & E4 (portfolio $-$ \texttt{best\_of\_n}) & +14.7 & [+8.4, +21.0] & $<10^{-4}$ \\
\midrule
\textsc{deepseek-chat} & E5 (\texttt{vf} $-$ \texttt{vf\_masked}) & +0.0 & [-4.6, +4.6] & 1.0000 \\
\textsc{qwen-plus} & E5 (\texttt{vf} $-$ \texttt{vf\_masked}) & +6.4 & [+2.2, +10.7] & 0.0065 \\
\textsc{deepseek-chat} & E3 (\texttt{vf}@1 $-$ \texttt{gr}@1) & +0.5 & [-3.4, +4.3] & 0.814$^\dagger$ \\
\textsc{qwen-plus} & E3 (\texttt{vf}@1 $-$ \texttt{gr}@1) & +1.8 & [-1.1, +4.8] & 0.226$^\dagger$ \\
\midrule
\multicolumn{5}{l}{\emph{preregistered \textsc{MiniMax-M2.5} extension (separate three-hypothesis Holm family)}} \\
\midrule
\textsc{MiniMax-M2.5} & E2 (\texttt{gr} $-$ \texttt{ss}) & +11.5 & [+5.9, +17.1] & 0.0002 \\
\textsc{MiniMax-M2.5} & E1 (\texttt{vf} $-$ \texttt{gr}) & +2.8 & [-2.0, +7.5] & 0.2590 \\
\textsc{MiniMax-M2.5} & E4 (portfolio $-$ \texttt{best\_of\_n}) & -9.0 & [-15.2, -2.7] & 0.0094 \\
\textsc{MiniMax-M2.5} & E5 (\texttt{vf} $-$ \texttt{vf\_masked}) & -0.7 & [-5.3, +3.9] & 0.770$^\dagger$ \\
\bottomrule
\end{tabular}

\caption{Frozen GEE endpoints (standardized marginal risk differences in
percentage points, fixture-clustered robust 95\% CIs). The primary family
\{E1, E2, E4\} $\times$ \{\dscode{}, \qwcode{}\} carries one paper-wide
Holm adjustment; E5 has its own two-hypothesis family;
$\dagger$: E3 is reported descriptively (unadjusted).}
\label{tab:primary}
\end{table}

\begin{table}[t]
\centering
\begin{tabular}{llccc}
\toprule
model & policy & rep 1 & rep 2 & rep 3 \\
\midrule
\textsc{deepseek-chat} & \texttt{ss} & 8/145 & 11/145 & 17/145 \\
 & \texttt{gr} & 22/145 & 24/145 & 26/145 \\
 & \texttt{vf} & 20/145 & 25/145 & 19/145 \\
 & \texttt{vf\_masked} & 15/145 & 25/145 & 24/145 \\
 & \texttt{best\_of\_n} & 58/145 & 65/145 & 59/145 \\
\midrule
\textsc{qwen-plus} & \texttt{ss} & 10/145 & 11/145 & 9/145 \\
 & \texttt{gr} & 19/145 & 18/145 & 24/145 \\
 & \texttt{vf} & 28/145 & 32/145 & 39/145 \\
 & \texttt{vf\_masked} & 13/145 & 30/145 & 28/145 \\
 & \texttt{best\_of\_n} & 21/145 & 26/145 & 27/145 \\
\bottomrule
\end{tabular}

\caption{Per-replication verifier-certified repair counts (out of 145)
for every policy, model, and replication.}
\label{tab:perrep}
\end{table}

\todo{prose for the three claims: (1) E2 replicated on both models;
(2) opposite-signed E1/E4/E5 => sampler-type (deepseek) vs
feedback-follower (qwen); note deepseek E5 RD exactly 0.0 (64 vs 64);
(3) certificate-as-invariant framing. Report validity-discordance
decomposition and per-class descriptive tables (appendix) with NO
significance language. Budget: Table~\ref{tab:cost}.}

\begin{table}[t]
\centering
\begin{tabular}{lrrr}
\toprule
model & model calls & input tokens & output tokens \\
\midrule
\textsc{deepseek-chat} & 9,139 & 17.8M & 6.6M \\
\textsc{qwen-plus} & 8,944 & 14.7M & 6.6M \\
\bottomrule
\end{tabular}

\caption{Confirmatory-campaign compute accounting (all calls including
failures; protocol section 11).}
\label{tab:cost}
\end{table}

\section{Exploratory extensions}
\todo{depth-2 collapse (difficulty dial; all strategies to the floor);
frustration cheating (copy-rate explodes with difficulty, guard
integrity); interface floor across model ladder (gemini-lite/8B below
floor; frontier-lite pilot: gemini-3.5-flash 0/17 single-shot; Sonnet
4.6 blind-detection degeneracy); per-class heterogeneity descriptive
forest plot; MiniMax scout/extension IF the scout is clean
\todo{pending scout}.}

\section{Related work}
\todo{AI4Math verifier-in-the-loop (AlphaProof, LeanDojo lineage);
LLM self-correction literature (intrinsic self-correction fails;
external verifiers help); rejection sampling / best-of-N with reward
models vs our exact symbolic filter; physics-domain LLM benchmarks.}

\section{Limitations and outlook}
\todo{two models => existence claim only, no strategy typology;
single benchmark domain (4d N=1 pure-quiver Seiberg); certificates not
proofs; depth-2 unsolved by all tested strategies; frontier suite not
run (cost symmetry argument + pilot exhibits); masked-feedback
placeholder-attention caveat (protocol wording verbatim).}

\appendix
\section{Frozen analysis protocol}
\todo{include paper/analysis\_protocol.md verbatim or as artifact
pointer + freeze/manifest commit hashes 813fdb0 / c155cff.}
\section{Reproduction}
\todo{commands, seeds, config hashes, expected costs; artifact layout.}

\bibliographystyle{unsrtnat}
\bibliography{refs}

\end{document}
